\begin{lemma}[Proper colouring of the reduced square]
Let $R$ be a finite simple graph of maximum degree at most $D$.  There are
$q\leq D^2+1$ nonempty available colours and a map
$c:V(R)\to[q]$ such that distinct vertices at distance at most two in $R$
receive different colours.  Consequently, vertices in one colour class are
nonadjacent and have pairwise disjoint $R$-neighbourhoods.
\end{lemma}
\begin{proof}
Form the square graph $S$ on $V(R)$ by joining two distinct vertices when
they are adjacent in $R$ or have a common neighbour in $R$.  Fix $v$.
There are at most $D$ neighbours of $v$ in $R$.  For each such neighbour
$u$, there are at most $D-1$ vertices other than $v$ adjacent to $u$.
Every neighbour of $v$ in $S$ is either one of the first set or belongs to
one of the latter sets.  Hence
$\deg_S(v)\leq D+D(D-1)=D^2$.  (Repeated vertices only decrease the size of
the union.)

Order the finite vertex set arbitrarily and colour it greedily.  When a
vertex is reached, at most $D^2$ colours have appeared on its already
coloured $S$-neighbours, so one of $D^2+1$ colours remains available.
Taking $q=D^2+1$ (which is positive) proves the asserted map and bound.
If two same-coloured vertices were adjacent in $R$, they would be adjacent
in $S$; if their $R$-neighbourhoods met, their common neighbour would again
make them adjacent in $S$.  Both contradict properness, proving the final
consequences.
\end{proof}
